Bei Otto von Guerickes Demonstration mit den Magdeburger Halbkugeln zogen auf jeder Seite \nPfO\ Pferde. Die Zugkraft eines Pferdes beträgt rund \FPfO. Schätzen Sie die Mindestgrösse der Halbkugeln ab.

Guericke selbst schreibt dazu in seinem Hauptwerk \emph{Experimenta nova (ut vocantur) Magdeburgica De Vacuo Spatio}: „Mit dem Lederring als Zwischenlage wurden nun diese Halbkugeln aufeinandergepaßt und dann die Luft \dots\ rasch ausgepumpt. Da sah ich, mit wieviel Gewalt sich die beiden Schalen gegen den Ring preßten! Und diesergestalt hafteten sie unter der Einwirkung des Luftdrucks so fest aneinander, daß 16 Pferde sie gar nicht oder nur mühsam auseinanderzureißen vermochten. Gelingt aber bei größter Kraftanstrengung die Trennung zuweilen doch noch, so gibt es einen Knall wie von einem Büchsenschuß.“
